% ============================================================
%  Nogor Shomadhan -- Report Sections 1 and 5
%  Engine: pdfLaTeX
% ============================================================
\documentclass[11pt,a4paper]{article}

% ---------- Encoding & fonts (pdfLaTeX-safe) ----------
\usepackage[T1]{fontenc}
\usepackage[utf8]{inputenc}
\usepackage{lmodern}
\usepackage{microtype}

% ---------- Layout ----------
\usepackage[a4paper,margin=1in]{geometry}
\usepackage{setspace}
\onehalfspacing

% ---------- Tables ----------
\usepackage{booktabs}
\usepackage{tabularx}
\usepackage{longtable}
\usepackage{xltabular}   % multi-page table with an X (auto-width) column
\usepackage{array}
\usepackage{ragged2e}
\newcolumntype{L}[1]{>{\RaggedRight\arraybackslash}p{#1}}

% ---------- Lists, colour, links ----------
\usepackage{enumitem}
\usepackage{xcolor}
\usepackage[hidelinks]{hyperref}
\usepackage{graphicx}

% ---------- Code / identifier styling ----------
\newcommand{\code}[1]{\texttt{\small #1}}

\begin{document}

% ============================================================
\section{Introduction}
% ============================================================

\subsection{Project Title}

\textbf{Nogor Shomadhan} --- \emph{A Mobile Platform for Community Maintenance
Issue Reporting, AI-Assisted Triage and Resolution Tracking}.

The name \emph{Nogor Shomadhan} translates from Bengali as ``City Solution'',
reflecting the aim of the project: turning scattered, informal complaints about a
neighbourhood into an organised, traceable and measurable resolution process.

\subsection{Project Overview}

Residents of a residential community routinely face local infrastructure and
maintenance problems --- damaged roads, faulty streetlights, water leakage,
overflowing garbage bins, blocked drainage, broken park facilities and
security-related concerns. In most communities these problems are still reported
through phone calls, social media groups or informal messaging, so complaints
arrive through disconnected channels, duplicate reports pile up, evidence is
missing, and residents have no way to tell whether anything is being done.

\emph{Nogor Shomadhan} replaces that informal process with a single mobile
application that carries a complaint through its entire life cycle. A resident
submits a report containing a title, a written description, a structured location
(house, road, avenue and nearby landmark) and photographic evidence. A
generative-AI model then reads the submitted text together with the attached
photographs and automatically assigns the complaint to one of twelve civic
categories, so the resident is never asked to classify the problem manually. The
same AI layer immediately compares the new report against recent complaints in
the same category and, when it recognises the same real-world problem, raises a
duplicate flag together with a confidence score and a written justification for a
human reviewer, instead of silently discarding the report.

An administrator verifies the complaint, resolves any duplicate flag, and
forwards genuine complaints to the community authority. The authority manages the
complaint from a dedicated dashboard: it accepts the work, assigns a contractor,
sets and revises deadlines, posts progress updates with photographic evidence,
and finally marks the complaint resolved. Every status transition is recorded in
an immutable history, so the resident sees a complete, timestamped timeline of
what happened rather than a single opaque status label. Once the issue is closed
the resident rates the quality of the resolution and can leave written feedback,
which the authority answers from a feedback centre. Any resident may also raise
the urgency counter on an existing complaint, letting the community collectively
signal which problems need attention first.

Around this core workflow the platform provides a role-aware notification system
that is generated inside the database itself, so residents, authority staff and
administrators are informed of every relevant event; a community forum where
authorities publish announcements, updates and alerts and residents hold
moderated discussions; analytics dashboards that visualise complaint distribution
by category and by area, monthly trends, resolved-versus-pending ratios and
resolution-time statistics; exportable PDF reports generated on the device; and
an administrative activity log and system-settings panel that make platform-level
actions auditable and configurable.

\subsection{Target Users}

The platform is built around three primary, clearly separated user roles. Each
role is a distinct route group in the application with its own navigation,
permissions and dashboard, so a user only ever sees the screens and data relevant
to their responsibility.

\begin{table}[htbp]
\centering
\caption{Target users of the Nogor Shomadhan platform}
\label{tab:target-users}
\renewcommand{\arraystretch}{1.25}
\begin{tabularx}{\textwidth}{L{2.6cm} L{3.4cm} X}
\toprule
\textbf{User Role} & \textbf{Who They Are} & \textbf{What They Do on the Platform} \\
\midrule
Resident &
Verified residents of the community (flat owners and tenants) &
Register an account for approval; submit complaints with photographs, a
structured location and a description; track the status timeline of their own
reports; browse and comment on the community complaint feed; raise the urgency
counter on existing complaints; rate and give feedback on resolved issues;
participate in the community forum; view personal analytics and export PDF
reports. \\
\addlinespace
Community Authority &
Staff of the community management or maintenance office &
Review complaints forwarded by the administrator; accept or reject work; assign
contractors; set and revise deadlines; post progress updates with evidence
photographs; move complaints through the \emph{Pending} $\rightarrow$
\emph{In Progress} $\rightarrow$ \emph{Resolved} life cycle; respond to resident
feedback from the feedback centre; publish official announcements, updates and
alerts to the forum; monitor operational analytics and resolution performance. \\
\addlinespace
Administrator &
Platform owner or community management committee representative &
Approve or reject resident and authority account registrations; verify submitted
complaints before they reach the authority; adjudicate AI-flagged duplicate
reports; moderate forum content; configure system settings such as AI
auto-categorisation, duplicate-detection sensitivity and maintenance mode; review
the platform-wide activity log and system analytics. \\
\bottomrule
\end{tabularx}
\end{table}

Beyond these three operational roles, the platform also serves the community
management committee as a secondary, indirect audience: the analytics dashboards
and exportable PDF reports turn the accumulated complaint history into evidence
for budgeting, contractor evaluation and long-term maintenance planning. Finally,
a set of public, unauthenticated pages (Home, About, Impact and Contact)
addresses prospective residents and visitors, communicating what the platform
does before they register.

% ------------------------------------------------------------
% Sections 2, 3 and 4 are written by other team members.
% The counter below makes the next section number 5.
% Delete this line once all sections are merged in order.
% ------------------------------------------------------------
\setcounter{section}{4}

% ============================================================
\section{Technologies and Integrations}
% ============================================================

This section documents the frameworks, external services and AI models used to
build \emph{Nogor Shomadhan}, and explains why each was selected over the
alternatives that were considered. The guiding constraints were that the system
had to be delivered by a small student team within a single semester, run on real
Android devices as well as in a browser for demonstration, and support a
relational, audit-heavy data model without the team having to build, deploy and
maintain a separate backend server.

\subsection{Technology Stack Overview}

Table~\ref{tab:stack} lists the technologies that make up the platform, grouped
by the layer of the system they serve. Each of them is described, and its
selection justified, in the subsections that follow.

\begin{table}[!htbp]
\centering
\caption{Summary of the technology stack}
\label{tab:stack}
\renewcommand{\arraystretch}{1.25}
\begin{tabularx}{\textwidth}{L{3.3cm} L{4.1cm} X}
\toprule
\textbf{Layer} & \textbf{Technology} & \textbf{Role in the Project} \\
\midrule
Application framework & Expo SDK~57 with React Native~0.86 and React~19 &
Cross-platform mobile application for Android and iOS, plus a web build for
demonstration \\
\addlinespace
Language & TypeScript~6 (\code{strict} mode) &
Statically typed application code, with shared domain types derived from the ER
diagram \\
\addlinespace
Navigation & Expo Router~57 (file-based routing, typed routes) &
Route groups per user role with per-group layout guards \\
\addlinespace
Backend platform & Supabase (PostgreSQL, Storage, PL/pgSQL, \code{pg\_cron}) &
Database, file hosting, business-rule triggers and scheduled jobs \\
\addlinespace
Data access & \code{@supabase/supabase-js} v2 &
Typed client-side queries, filtering and mutations against PostgreSQL \\
\addlinespace
Artificial intelligence & Google Gemini Flash via \code{@google/generative-ai} &
Multimodal complaint categorisation and duplicate-report detection \\
\addlinespace
Messaging integration & BulkSMSBD HTTP SMS API &
Out-of-band SMS notices for account approval and rejection \\
\addlinespace
Device capabilities & \code{expo-image-picker}, \code{expo-image},
\code{expo-file-system}, \code{expo-print}, \code{expo-sharing},
\code{expo-secure-store} &
Evidence capture, image rendering, on-device PDF report generation and sharing \\
\addlinespace
UI and interaction & React Native Reanimated~4, Gesture Handler, Safe Area
Context, Keyboard Controller, \code{@expo/vector-icons} &
Animated, gesture-driven and accessible interface components \\
\addlinespace
Local persistence & \code{@react-native-async-storage/\allowbreak async-storage} &
Session persistence, cached notification state and offline-tolerant preferences \\
\addlinespace
Tooling & Expo CLI, Metro bundler, ESLint (\code{eslint-config-expo}), Git &
Development server, bundling, static analysis and version control \\
\bottomrule
\end{tabularx}
\end{table}

\subsection{Frontend Framework: Expo and React Native}

The application is built with \textbf{Expo SDK~57} on top of
\textbf{React Native~0.86} and \textbf{React~19}, written entirely in
\textbf{TypeScript} with \code{strict} type checking enabled. Routing is handled
by \textbf{Expo Router~57}, which maps the directory structure under
\code{src/app/} directly onto application routes.

Expo Router's route groups are used to express the role architecture of the
system: \code{(public)}, \code{(resident)}, \code{(authority)} and \code{(admin)}
each form an isolated navigation tree with its own \code{\_layout.tsx} file that
acts as an access guard for that role. A resident can therefore never reach an
administrator screen through navigation, because the guard lives at the layout
level rather than being repeated inside every screen. Typed routes were enabled
so that an incorrect navigation target becomes a compile-time error instead of a
runtime dead end.

\paragraph{Justification.} Three alternatives were considered for the client
application:

\begin{itemize}[leftmargin=1.4em,itemsep=0.35em]
  \item \textbf{Native Android (Kotlin) and iOS (Swift) development} would have
        required two separate codebases, two toolchains and two sets of UI work
        for the same feature set --- not feasible for a four-member team within
        one semester.
  \item \textbf{Flutter} offers comparable cross-platform reach, but would have
        required the team to adopt Dart and would have cut the project off from
        the JavaScript ecosystem that the Supabase and Google Generative AI
        client libraries are written for. The team's existing competence in
        JavaScript and TypeScript made React Native the lower-risk choice.
  \item \textbf{Bare React Native} (React Native CLI without Expo) would have
        given the same rendering model, but the project depends on more than
        twenty native modules --- camera and gallery access, secure storage, file
        system, printing, sharing, gesture handling, animation worklets. Under
        bare React Native each of these requires manual native linking and
        per-platform build configuration. Expo's config plugins and prebuild
        pipeline handle that automatically, which removed a large and unforgiving
        class of build problems from the project.
\end{itemize}

Expo was additionally preferred because \textbf{React Native Web} is part of the
same toolchain: the identical codebase produces a browser build, which made
in-class demonstration and day-to-day development markedly faster than launching
an emulator for every change. Expo Router was chosen over configuring
\textbf{React Navigation} manually because file-based routing keeps the
navigation structure visible in the file tree and makes role-based guarding a
structural property of the project rather than a convention the team has to
remember to apply.

\subsection{Backend and Database Platform: Supabase}

\textbf{Supabase} provides the entire server side of the platform. Four of its
capabilities are used.

\begin{description}[leftmargin=1.6em,style=nextline,itemsep=0.4em]
  \item[PostgreSQL database.] The relational schema implements the ER diagram
        directly: \code{account}, \code{complaints}, \code{evidence},
        \code{duplicate}, \code{complaint\_status\_history},
        \code{complaint\_work\_updates}, \code{contractor\_history}, feedback,
        forum posts and comments, \code{notifications} and \code{app\_settings}.
        Domain vocabularies such as account status, account role, complaint
        status, complaint category and notification type are declared as
        PostgreSQL \code{ENUM} types, so invalid values are rejected by the
        database rather than merely discouraged by the client. Row Level Security
        is enabled on every table with explicit access policies.
  \item[Storage.] Evidence photographs submitted by residents and progress
        photographs uploaded by the authority are stored in Supabase Storage
        buckets. Images captured through the picker are decoded from base64 into
        a binary buffer before upload, and the resulting public URL is written
        into the \code{evidence} table alongside the complaint.
  \item[PL/pgSQL functions and triggers.] The notification system is implemented
        inside the database. Triggers on account insertion, complaint status
        changes, duplicate insertion, duplicate adjudication, work updates,
        feedback, feedback replies, forum comments and official announcements
        call a shared \code{enqueue\_notification} routine that writes
        role-addressed notification rows carrying a deduplicating event key, a
        priority level and a deep-link action path.
  \item[Scheduled jobs.] The \code{pg\_cron} extension runs time-based
        notification generators inside the database --- for example, flagging a
        complaint that has remained pending for fourteen days without work being
        started, and raising deadline-milestone and overdue notices.
\end{description}

\paragraph{Justification.} Supabase was chosen over the two obvious
alternatives:

\begin{itemize}[leftmargin=1.4em,itemsep=0.35em]
  \item \textbf{A custom backend (Node.js/Express with PostgreSQL)} would have
        given complete control, but the team would then have had to design,
        implement, document, deploy and keep a REST API online in addition to
        building the mobile client. Supabase exposes PostgreSQL through a
        generated, policy-enforced client SDK, so the endpoints planned in the
        API section of this report are satisfied by typed queries from the
        application's service layer, with no server code to host. This removed an
        entire deployment surface from a semester-length project.
  \item \textbf{Firebase (Firestore)} was rejected because the data model is
        strongly relational. The platform constantly joins complaints to
        accounts, evidence, status history, contractor history, feedback and
        duplicate records, and the analytics screens aggregate across those joins
        by category, area, month and resolution time. A document store would have
        forced denormalisation and client-side aggregation, and it offers no
        equivalent of \code{ENUM} domains, foreign keys or transactional
        triggers. PostgreSQL expresses the ER diagram as designed.
\end{itemize}

Implementing notifications as database triggers rather than as client-side code
or separate cloud functions was a deliberate decision: a notification row is
created in the same transaction as the state change that caused it, so
notifications cannot drift out of step with the data, and they are produced
correctly no matter which role or which screen triggered the change. For the same
reason \code{pg\_cron} was preferred over an external scheduler or worker process
--- the schedule lives next to the data it inspects, with no additional service
to keep running.

\subsection{AI Models and APIs}

Artificial intelligence is used for the two smart features of the platform, both
served by \textbf{Google Gemini Flash} through the official
\textbf{\code{@google/generative-ai}} JavaScript SDK.

\subsubsection{Automatic Complaint Categorisation}

When a resident submits a complaint, the title, the description and the attached
evidence photographs are sent to Gemini in a single multimodal request. The model
is constrained to reply with exactly one label from a fixed list of twelve
categories --- Road Damage; Garbage \& Waste; Drainage \& Waterlogging;
Streetlight \& Electrical; Water Supply; Sanitation \& Public Toilets; Traffic \&
Illegal Parking; Public Safety \& Encroachment; Noise \& Environmental Pollution;
Parks \& Public Spaces; Animal-Related Issues; and Other. The returned label is
validated against that list before it is accepted, and the system falls back to
\emph{Other} if the response is unrecognised or the service is unavailable, so
complaint submission never fails because of the AI layer.

\subsubsection{Duplicate Complaint Detection}

Duplicate detection runs as a two-stage pipeline:

\begin{enumerate}[leftmargin=1.6em,itemsep=0.35em]
  \item \textbf{Database prefilter.} A PostgreSQL query retrieves a small
        candidate set --- the most recent complaints in the same category,
        excluding the new complaint itself. This keeps the AI request small and
        cheap regardless of how large the complaint history grows.
  \item \textbf{AI adjudication.} The new complaint and the candidate set are
        passed to Gemini, which is instructed to reason about physical location
        (house, road, avenue, landmark), category and timing rather than mere word
        overlap, and to return a strict JSON verdict containing a boolean
        duplicate decision, a confidence score, the matched complaint identifier
        and a written reason. When the score crosses the configured threshold, a
        row is inserted into the \code{duplicate} table carrying the score and the
        model's reasoning, which in turn raises a duplicate-review notification
        for the administrator.
\end{enumerate}

Crucially, the AI never deletes a complaint. It produces evidence; an
administrator makes the decision on a dedicated review screen. The sensitivity of
the check is configurable by the administrator from the system-settings page, and
both AI features can be switched off entirely through the \code{app\_settings}
table.

\paragraph{Justification.} Several approaches were evaluated:

\begin{itemize}[leftmargin=1.4em,itemsep=0.35em]
  \item \textbf{Training a custom classifier} (for example a fine-tuned BERT
        model or a classical TF--IDF classifier) would have required a labelled
        corpus of Bangladeshi community complaints. No such dataset exists, and
        collecting and annotating one was outside the scope of the project. A
        general-purpose instruction-following model achieves usable accuracy with
        no training data at all.
  \item \textbf{Keyword matching or TF--IDF cosine similarity} for duplicate
        detection was rejected because the reports it must compare are written by
        different residents in free text, frequently mixing Bengali and English
        and describing the same pothole in entirely different words. A lexical
        measure misses paraphrase, and it cannot weigh location and recency the
        way the task requires.
  \item \textbf{Vector embeddings with \code{pgvector}} were a serious candidate
        and would scale better on a very large corpus. They were not chosen
        because a similarity score alone cannot explain itself: the administrator
        review screen needs a human-readable justification for why two complaints
        were judged to be the same problem. A single generative call returns the
        score \emph{and} the reasoning, which made the review workflow possible
        without building a second explanation mechanism.
  \item \textbf{Other commercial models} (such as OpenAI GPT-4 class or Anthropic
        Claude models) are comparable in capability. Gemini Flash was selected
        because its free tier is sufficient for a student project, it accepts text
        and images in the same request --- which matters because the photograph
        often identifies the category better than the resident's wording --- and
        its latency is low enough to run inside the complaint submission flow
        without making the user wait noticeably.
\end{itemize}

\subsection{External APIs and Third-Party Services}

\paragraph{BulkSMSBD SMS gateway.} Account approval and rejection decisions are
delivered to the applicant by SMS through the BulkSMSBD HTTP API. This closes a
genuine gap in the workflow: a resident whose registration is still awaiting
approval cannot sign in, and therefore cannot receive an in-app notification. A
channel outside the application is the only way to reach them. BulkSMSBD was
chosen over international providers such as Twilio because it supports
Bangladeshi eleven-digit mobile numbers natively with local operator routing, its
pricing suits a local deployment, and it requires no international sender
registration. Email was considered as a cheaper alternative but rejected because
SMS reach and read rates are substantially higher among the intended user base.

\paragraph{On-device PDF report generation.} Analytics reports and complaint
reports are produced with \code{expo-print}, which renders an HTML document
generated by the application into a PDF, and \code{expo-sharing}, which hands the
resulting file to the operating system share sheet. Server-side generation with a
headless browser toolchain such as Puppeteer or \code{wkhtmltopdf} was rejected
because it would have meant hosting yet another service purely to produce a
document the device can already render, and because on-device generation keeps
complaint data on the user's device.

\paragraph{Device and media integrations.} Evidence capture uses
\code{expo-image-picker} to take a photograph or select one from the gallery,
returning base64 data that is uploaded to Supabase Storage;
\code{expo-file-system} manages generated report files; \code{expo-image}
provides cached, progressive image rendering; \code{expo-secure-store} holds
sensitive session material on the device; and
\code{@react-native-async-storage/async-storage} persists the Supabase session
and cached notification state so the user is not signed out between launches.
Interface behaviour relies on \textbf{React Native Reanimated~4} with the
worklets runtime for animations that run on the UI thread, \textbf{React Native
Gesture Handler} for native-quality gestures, \textbf{React Native Safe Area
Context} and \textbf{Screens} for correct native layout, and \textbf{React Native
Keyboard Controller} to keep long complaint and registration forms usable while
the keyboard is open. The analytics charts are composed from React Native
primitives and the project's own theme tokens rather than from a heavyweight
charting library, which kept the visualisations consistent with the rest of the
interface and avoided a dependency that renders differently on web and native.

\paragraph{State management.} Application state is managed with React Context
providers and custom hooks (for authentication, complaints, notifications and
analytics) rather than an external state library such as Redux or Zustand.
Because Supabase is the single source of truth and each screen queries the data it
needs, a global normalised client store would have added indirection without
solving a problem the project actually has.

\subsection{Development Tooling}

Development uses the \textbf{Expo CLI} with the \textbf{Metro} bundler for the
development server and bundling, \textbf{ESLint} with \code{eslint-config-expo}
for static analysis, \textbf{TypeScript} path aliases (\code{@/*}
$\rightarrow$ \code{src/*}) to keep imports stable as the file tree grows, and
\textbf{Git} for version control. Database changes are tracked as versioned SQL
migration files inside the repository, and the current schema and Row Level
Security policies are maintained as living documentation alongside them, so the
database can be rebuilt from scratch in a new Supabase project.

\subsection{Consolidated Justification}

Table~\ref{tab:justification} summarises, for each major requirement of the
system, the technology adopted, the alternatives that were considered, and the
reason for the final decision.

\begin{xltabular}{\textwidth}{L{2.7cm} L{2.9cm} L{2.9cm} X}
\caption{Technology decisions and their justification}
\label{tab:justification} \\
\toprule
\textbf{Requirement} & \textbf{Adopted} & \textbf{Alternatives} &
\textbf{Reason for Choice} \\
\midrule
\endfirsthead
\toprule
\textbf{Requirement} & \textbf{Adopted} & \textbf{Alternatives} &
\textbf{Reason for Choice} \\
\midrule
\endhead
\bottomrule
\endfoot
\bottomrule
\endlastfoot

Cross-platform mobile client &
Expo (React Native) &
Native Kotlin/Swift; Flutter; bare React Native &
One TypeScript codebase for Android, iOS and web; managed native modules remove
manual linking for twenty-plus native dependencies; team already fluent in
TypeScript \\
\addlinespace

Role-separated navigation &
Expo Router with route groups &
Manually configured React Navigation &
File-based routes mirror the four user modules; role guards live in per-group
layouts; typed routes catch invalid navigation at compile time \\
\addlinespace

Relational data with audit history &
Supabase PostgreSQL &
Firebase Firestore; self-hosted Express and PostgreSQL &
Foreign keys, \code{ENUM} domains and transactional triggers express the ER
diagram directly; the generated client SDK removes the need to host a separate
API \\
\addlinespace

Evidence image hosting &
Supabase Storage &
Cloudinary; Amazon S3 &
Same project, credentials and access-policy model as the database; public URLs
are referenced straight from complaint rows \\
\addlinespace

Event-driven notifications &
PL/pgSQL triggers &
Client-side writes; serverless cloud functions &
Notifications are written in the same transaction as the state change, so they
can never drift from the data, whichever client caused the change \\
\addlinespace

Time-based notifications &
\code{pg\_cron} &
External cron host; scheduled worker service &
Runs next to the data it inspects, with no extra service to deploy or monitor \\
\addlinespace

Complaint categorisation &
Google Gemini Flash (multimodal) &
Fine-tuned BERT; TF--IDF classifier; other commercial LLMs &
No labelled Bangladeshi complaint corpus exists; photographs and text are
classified in a single request; free tier and low latency fit the submission
flow \\
\addlinespace

Duplicate detection &
Gemini with an SQL candidate prefilter &
Keyword or TF--IDF similarity; \code{pgvector} embeddings &
Recognises semantically identical reports written in different words and
languages, weighs location and recency, and returns a human-readable reason for
administrator review \\
\addlinespace

Out-of-band user notices &
BulkSMSBD API &
Twilio; transactional email; push notifications &
Local eleven-digit number support and operator routing at local cost; reaches
users who cannot yet sign in to receive an in-app notification \\
\addlinespace

Report export &
\code{expo-print} with \code{expo-sharing} &
Server-side Puppeteer or \code{wkhtmltopdf} &
Renders PDFs on the device, so no additional service is hosted and complaint data
stays local \\
\addlinespace

Client state &
React Context with custom hooks &
Redux Toolkit; Zustand &
Supabase is the single source of truth; per-screen queries make a global
normalised store unnecessary \\

\end{xltabular}

\end{document}
